\chap{1}{INTRODUCTION}{-1pt}
Genetic variant interpretation is the task of deciding whether a change in a person's DNA is harmless or disease-causing. A single base change in a gene such as \textit{BRCA1} can leave the protein untouched, or it can switch the gene off and sharply raise the lifetime risk of breast and ovarian cancer. Telling the two apart is one of the central problems of clinical genetics. Laboratories classify variants with the ACMG/AMP guidelines, which combine population frequency, family segregation, functional assays and computational predictions into one of five classes \cite{richards2015acmg}. The process is reliable, but it is manual, slow, and least effective for rare variants that have never been observed before.

The \emph{genome} is the complete set of DNA of an organism. DNA is a long molecule written in a four-letter alphabet of nucleotides, A, C, G and T, and the human genome contains about 3.1 billion base pairs, divided among 22 numbered chromosomes, the sex chromosomes X and Y and a small mitochondrial genome. Along these chromosomes lie about twenty thousand protein-coding genes. Each gene is made of exons, which carry the message for the protein, and introns, which are cut out, and its activity is steered by regulatory DNA nearby. Only one to two per cent of the genome codes for protein, so a change in DNA can matter both inside and outside the exons.

Any two people differ from each other, and from the reference, at several million positions. Most of these \emph{variants} are harmless, but some change a protein or disturb the control of a gene. They range from a single-nucleotide variant (SNV), the replacement of one base by another, to insertions, deletions and combined deletion--insertions of longer stretches. To say where a variant lies, scientists number every position of a \emph{reference genome}, a single agreed sequence. The human reference exists in several versions called assemblies, of which GRCh37 (hg19) and GRCh38 (hg38) are in common use, and the same base has a different coordinate in each. A position is therefore meaningful only together with the name of its assembly, a fact that shapes the design of Evo2Serve.

The gene used throughout this project, \textit{BRCA1}, lies on the long arm of chromosome 17 (17q21.31) and spans about 126 thousand bases in GRCh38. Its protein takes part in the repair of double-stranded breaks in DNA, and inherited loss-of-function variants raise the lifetime risk of breast and ovarian cancer. Because thousands of single-base changes in this gene have been tested experimentally \cite{findlay2018brca1}, it is the best-characterised benchmark for variant interpretation.

As genetic testing becomes cheaper, laboratories report more and more \textit{variants of uncertain significance}, that is, variants for which the available evidence is insufficient to reach a verdict. Fast computational tools that give a first, model-based opinion on such variants are therefore valuable, both to researchers who must prioritise variants for laboratory follow-up and to students who are learning how variant interpretation works. Classical predictors such as SIFT, PolyPhen-2 and CADD rely on hand-crafted features, multiple sequence alignments and precomputed annotations \cite{ng2001sift,adzhubei2010polyphen2,rentzsch2019cadd}. They work well for protein-coding changes, but many are restricted to missense substitutions and none of them can be asked about an arbitrary, user-defined change in an arbitrary part of the genome.

Genomic language models offer a different route. A language model reads a sequence and learns the probability of every next symbol from the symbols before it. A DNA language model trained on enough genomes learns what ``natural'' DNA looks like, so a mutation that makes a sequence look less natural is likely to be disruptive. Evo~2, released by the Arc Institute, is the largest open model of this kind. It was trained on trillions of nucleotides from all domains of life, models DNA at single-nucleotide resolution with a context of up to one million bases, and, without any task-specific training, reaches an AUROC of 0.9037 on the \textit{BRCA1} saturation genome editing benchmark \cite{brixi2026evo2,findlay2018brca1}.

However, such a model is distributed as research code and Jupyter notebooks that need a data-centre GPU and programming expertise. A student, a counsellor or a small laboratory cannot simply open a browser, choose a gene, type a mutation and ask the model what it thinks. Between the model and the person who wants an answer sits a software layer: a coordinate must be resolved against a reference genome, a context window must be retrieved, the reference and the mutated sequences must be scored, and the score difference must be turned into a label and a confidence value. Each of these steps can fail silently, and a mistake in any of them produces a confident answer about the wrong base.
\par\noindent The work reported here builds that software layer and examines it critically. \textbf{Evo2Serve} is a full-stack web application that deploys Evo~2-7B on a serverless NVIDIA H100 GPU through Modal and exposes it through an authenticated FastAPI endpoint. A Next.js web interface lets a user select a human genome assembly, browse chromosomes or search for a gene such as \textit{BRCA1}, read the reference sequence, list the known ClinVar variants of the gene, enter a mutation of their own, and compare the model's prediction with the expert ClinVar classification. The interface and the model service are described in Chapters~3 to~5.

The model score is converted to a label with a threshold calibrated on the \textit{BRCA1} saturation genome editing data of Findlay \textit{et al.}, in which nearly four thousand single-nucleotide variants were each tested for their effect on cell survival \cite{findlay2018brca1}. Because the quality of the label depends entirely on the surrounding software, the project does not stop at building the tool: it also verifies the deployed pipeline. We recompute the composition of the calibration data, evaluate five established predictors on the same rows, check that the context window and the genome assembly used at calibration time are the ones used at prediction time, and run a generated suite of 55 request cases against both validation layers of the service. The findings are reported in Chapter~6 together with the improvements they suggest.

Evo2Serve is a research and teaching tool. Under the ACMG/AMP framework a computational prediction is only supporting evidence and must not be the sole basis for a clinical assertion \cite{richards2015acmg}, and every label produced by the system is accordingly worded as ``likely'' and flagged as provisional where the calibration does not apply. The remainder of this chapter states the problem, the aim and objectives of the project, and its scope.
\label{ch:1:intro}

\vspace{6pt}
\sect{1.1}{PROBLEM STATEMENT}
\vspace{8pt}
Variant effect prediction has progressed rapidly, but the best models remain out of reach of the people who could benefit from them. State-of-the-art genomic language models such as Evo~2 are available only as research code that needs a high-end GPU, a specific software environment and programming skills. Public interpretation websites, on the other hand, display existing annotations and curated evidence; they do not let a user run a large genomic language model on a new mutation.

There is also a reliability problem. Turning a model score into an answer involves a chain of software steps: resolving the coordinate in the correct genome assembly, fetching the right context window, applying the mutation correctly, scoring two sequences, and comparing the difference with a threshold that was calibrated on specific data. A mistake at any step, for example a coordinate typed in the wrong assembly, produces a confident but meaningless answer, and the interface gives the user no warning.

The motivation behind developing the Evo2Serve system includes:

\begin{itemize}
\item \textbf{Accessibility:} A browser is enough to use a 7-billion-parameter DNA model; no GPU, installation or programming is needed.
\item \textbf{Speed:} A prediction is returned within seconds, compared with the weeks that expert curation of a rare variant can take.
\item \textbf{Coverage of all mutation types:} Substitutions, deletions, insertions and deletion--insertions are all supported, whereas many classical predictors handle only missense changes.
\item \textbf{Transparency:} Each prediction is shown next to the expert ClinVar classification, so users see where the model agrees and disagrees with curated evidence.
\item \textbf{Cost efficiency:} A serverless GPU is billed only while a request is processed and scales to zero when idle, so a college or small laboratory can host the service.
\item \textbf{Trustworthiness:} The deployed pipeline is verified against its calibration conditions, and its limits are stated openly.
\end{itemize}
\sectend{1.1}

\newpage
\sect{1.2}{AIM OF THE PROJECT}
To develop and verify an end-to-end, AI-based web application that deploys the Evo~2 genomic language model as a secure online service and uses it to analyse genetic mutations: a user selects a genome assembly, finds a gene, enters a single-nucleotide or insertion/deletion mutation, and receives within seconds a model-based pathogenicity label with a confidence value, shown beside the known ClinVar classification. The system combines public genome and variant databases, a serverless H100 GPU endpoint, authenticated access, a calibrated decision threshold and a rigorous evaluation of the deployed pipeline, so that the tool can serve education and research triage while stating clearly that its output is computational evidence and not a clinical diagnosis.
\sectend{1.2}

\sect{1.3}{OBJECTIVES OF THE PROJECT}
\begin{enumerate}
\item \textbf{Deploy zero-shot variant scoring as a secure service:} Serve Evo~2-7B on a serverless NVIDIA H100 GPU with Modal behind a FastAPI endpoint that accepts a single JSON request and is protected by signed session tokens and two validation layers. Score the reference and the mutated 8\,kb windows and compute the log-likelihood difference $\Delta$ for substitutions, deletions, insertions and deletion--insertions of up to 500 bases.
\item \textbf{Calibrate and report the decision:} Gather reference sequences, gene and variant records and the \textit{BRCA1} saturation genome editing workbook, select the threshold by maximising Youden's index on the ROC curve of $-\Delta$, and report a label, a confidence value and a provisional flag for every prediction.
\item \textbf{Build a genome browser interface with ClinVar comparison:} Let users choose a human assembly, browse chromosomes, search genes, read the reference sequence, list the known variants of a gene and see the Evo~2 prediction next to the expert ClinVar classification.
\item \textbf{Evaluate and verify the pipeline reproducibly:} Compare the model with CADD, PolyPhen-2, SIFT, Align-GVGD and phyloP on the same calibration rows, check the deployed window, assembly handling and request validation against the calibration conditions, provide scripts that reproduce every reported number without a GPU, and present all output as computational evidence, not as a diagnosis.
\end{enumerate}
\sectend{1.3}

\sect{1.4}{SCOPE OF THE PROJECT}
The scope of this project covers the complete design, implementation, deployment and verification of a full-stack genetic variant analysis platform. In the dataset view, the scope includes consuming reference sequences and gene coordinates from the UCSC Genome Browser, gene and variant records from NCBI and ClinVar, and the supplementary workbook of the \textit{BRCA1} saturation genome editing study. In the model view, it includes serving Evo~2-7B for zero-shot scoring of substitutions, deletions, insertions and deletion--insertions, with a decision threshold calibrated on \textit{BRCA1} substitutions. In the verification view, it includes the recomputation of the calibration data composition, the evaluation of five established predictors on the same rows, the check of the context window and assembly handling, and a generated suite of 55 request cases executed against both validation layers.

In the application view, the scope includes a responsive web interface with authentication, assembly selection, chromosome browsing, gene search, sequence display, a known-variant table, a variant entry form and a comparison view. The boundary of the project is deliberately drawn around the software layer. Retraining or fine-tuning Evo~2, wet-laboratory validation, clinical diagnosis and reporting, calibration for genes other than \textit{BRCA1}, and the measurement of latency and throughput on production traffic fall outside the scope, and are listed as future work.
\sectend{1.4}

\sect{1.5}{SYNOPSIS}
\vspace{4pt}
Evo2Serve is a web application that makes a state-of-the-art genomic language model usable from a browser for analysing genetic mutations and their likely disease impact. Interpreting a variant by hand is slow, and the best computational models are locked inside research notebooks. This project closes that gap by combining public genome and ClinVar data, a serverless H100 GPU running Evo~2-7B, token-based authentication and a calibrated decision rule into one platform that returns a label, a confidence value and a provisional flag within seconds.

Because a prediction is only as good as the coordinate and the sequence it is computed on, the project also verifies the pipeline it builds. It measures the calibration data, evaluates reference predictors on the same rows, checks the context window and the genome assembly, and tests the request validators, and it reports eight concrete improvements. The platform is intended for education and research triage and is explicit that its output is supporting computational evidence, not a diagnosis. The following chapters review the related literature, specify the requirements, describe the design and implementation, and report the results.
\sectend{1.5}
